\clearpage\pagestyle{otMain}\setlength{\otBlocEspacement}{2.5mm}\fontsize{11}{15}\selectfont\corrigesChapitre{1}{Les nombres jusqu’au milliard}
\begin{blocLivre}{cor-1-1}\begin{corrige}{1}{Écrire en chiffres}\demarche{Compléter chaque classe avec trois chiffres.}
On remplit classe par classe : $\res{3\,040\,200}$, $\res{60\,500\,000}$ et $\res{800\,000\,009}$.
\end{corrige}\end{blocLivre}
\begin{blocLivre}{cor-1-2}\begin{corrige}{2}{Lire par classes}\demarche{Lire chaque classe suivie de son nom.}
« $\res{\text{quarante-cinq-millions-deux-cent-huit-mille-trois-cent-seize}}$ ». La classe des mille porte $\res{208}$.
\end{corrige}\end{blocLivre}
\begin{blocLivre}{cor-1-3}\begin{corrige}{3}{La valeur des chiffres}\demarche{Repérer la classe avant le rang.}
Le chiffre des unités de millions est $\res5$. Le 3 occupe les centaines de millions : il vaut $\res{300\,000\,000}$.
\end{corrige}\end{blocLivre}
\begin{blocLivre}{cor-1-4}\begin{corrige}{4}{Nombre de milliers}\demarche{Séparer les trois derniers chiffres.}
Le chiffre des unités de mille est $\res8$. Le nombre de milliers est $\res{45\,208}$, car $45\,208\,316=45\,208\times1\,000+316$.
\end{corrige}\end{blocLivre}
\begin{blocLivre}{cor-1-5}\begin{corrige}{5}{Comparer}\demarche{Compter les chiffres puis comparer rang à rang.}
Sept chiffres contre huit : $\res{9\,987\,654<10\,000\,000}$. À longueur égale, les centaines de mille tranchent : $\res{45\,320\,000>45\,230\,000}$.
\end{corrige}\end{blocLivre}
\begin{blocLivre}{cor-1-6}\begin{corrige}{6}{Décomposer et recomposer}\demarche{N’écrire que les rangs non nuls.}
$7\,080\,205=\res{7\times1\,000\,000+8\times10\,000+2\times100+5}$. La recomposition donne $\res{60\,030\,500}$.
\end{corrige}\end{blocLivre}
\begin{blocLivre}{cor-1-7}\begin{corrige}{7}{Sur la demi-droite}\demarche{Chercher la valeur d’un petit intervalle.}
Deux graduations marquées voisines diffèrent de un million ; chaque petit pas vaut $200\,000$. On lit $\res{A(3\,200\,000)}$, $\res{B(4\,600\,000)}$ et $\res{C(5\,800\,000)}$.
\label{fig-cor-1-7}\figureGeo[corrige]{num-7}
\end{corrige}\end{blocLivre}
\begin{blocLivre}{cor-1-8}\begin{corrige}{8}{Encadrer}\demarche{Fixer le rang, puis prendre ses deux multiples voisins.}
$\res{7\,000\,000<7\,802\,540<8\,000\,000}$ et $\res{7\,800\,000<7\,802\,540<7\,810\,000}$.
\end{corrige}\end{blocLivre}
\begin{blocLivre}{cor-1-9}\begin{corrige}{9}{Intercaler}\demarche{Caler d’abord les classes, puis le rang imposé.}
Par exemple $\res{65\,000\,400}$ : il est entre les deux bornes et son chiffre des centaines est 4. D’autres choix conviennent.
\end{corrige}\end{blocLivre}
\begin{blocLivre}{cor-1-10}\begin{corrige}{10}{Ordonner}\demarche{Grouper par longueur avant de comparer.}
Les nombres à huit chiffres dominent : $\res{10\,100\,000>10\,010\,000>9\,900\,000>9\,090\,000}$.
\end{corrige}\end{blocLivre}
\begin{blocLivre}{cor-1-11}\begin{corrige}{11}{Les populations}\demarche{Comparer les dizaines de millions d’abord.}
L’ordre croissant est $\res{10\,500\,000<47\,400\,000<67\,800\,000<83\,200\,000}$. La plus grande vérifie $\res{80\,000\,000<83\,200\,000<90\,000\,000}$.
\end{corrige}\end{blocLivre}
\begin{blocLivre}{cor-1-12}\begin{corrige}{12}{Des paquets de mille}\demarche{Compter les milliers entiers.}
$2\,507\,000=2\,507\times1\,000$ : $\res{2\,507}$ cartons exactement. Avec 500 brochures de plus, un carton reste incomplet : toujours $\res{2\,507}$ cartons pleins et 500 brochures en attente.
\end{corrige}\end{blocLivre}
\begin{blocLivre}{cor-1-13}\begin{corrige}{13}{Le compteur du site}\demarche{S’appuyer sur le passage au rang supérieur.}
Il manque $10\,000\,000-9\,999\,500=\res{500}$ visites. Après 800 visites : $9\,999\,500+800=\res{10\,000\,300}$ : le cap des huit chiffres est franchi.
\end{corrige}\end{blocLivre}
\begin{blocLivre}{cor-1-14}\begin{corrige}{14}{Avec six cartes}\demarche{Choisir le chiffre de gauche en premier.}
Ordre décroissant : $\res{985\,210}$. Pour le plus petit, commencer par le plus petit chiffre non nul : $\res{102\,589}$.
\end{corrige}\end{blocLivre}
\begin{blocLivre}{cor-1-15}\begin{corrige}{15}{Vrai ou faux ?}\demarche{Compter les chiffres avant tout.}
On compare d’abord les longueurs : sept chiffres contre huit, donc $\res{9\,999\,999<10\,000\,001}$. Les chiffres utilisés ne décident pas. L’écart vaut $10\,000\,001-9\,999\,999=\res2$.
\end{corrige}\end{blocLivre}
